\begin{examplebox}
	\textbf{\textcolor{CExample}{例 1（基础）}}

	\textbf{\textcolor{CExample}{题目：}}
	圆 $O$ 的半径为 $5$，圆心到某弦的距离（弦心距）为 $3$，求该弦的长度。

	\textbf{\textcolor{CExample}{解题步骤：}}
	\begin{description}[style=nextline, leftmargin=2.8em, labelsep=0.5em, itemsep=0.45em]
		\item[\textbf{第一步（写公式）：}] 由弦长公式 $L=2\sqrt{r^{2}-d^{2}}$，代入 $r=5,d=3$。
			\par\textbf{理由：} 直接套用弦长公式。
		\item[\textbf{第二步（计算平方差）：}] $r^{2}-d^{2}=25-9=16$。
			\par\textbf{理由：} 数值计算。
		\item[\textbf{第三步（得结果）：}] $L=2\sqrt{16}=8$。
			\par\textbf{理由：} 开方后乘以 $2$ 得弦长。
	\end{description}

	\textbf{\textcolor{CExample}{关键结论：}}
	\textbf{弦长为 $8$。}

	\medskip
	\textbf{\textcolor{CExample}{例 2（稍变形）}}

	\textbf{\textcolor{CExample}{题目：}}
	圆 $C:(x-1)^{2}+(y+2)^{2}=25$ 被直线 $l:3x-4y+4=0$ 所截，求弦长。

	\textbf{\textcolor{CExample}{解题步骤：}}
	\begin{description}[style=nextline, leftmargin=2.8em, labelsep=0.5em, itemsep=0.45em]
		\item[\textbf{第一步（求弦心距）：}] 圆心 $C(1,-2)$，直线 $3x-4y+4=0$，距离
			\[
				\begin{aligned}
					d &= \frac{|3\cdot1-4\cdot(-2)+4|}{\sqrt{3^{2}+(-4)^{2}}} \\
					&= \frac{15}{5} \\
					&= 3
			\end{aligned}
			\]
			\par\textbf{理由：} 点到直线距离公式。
		\item[\textbf{第二步（确认半径）：}] 由圆方程知半径 $r=\sqrt{25}=5$。
			\par\textbf{理由：} 标准方程给出。
		\item[\textbf{第三步（代入公式）：}]
			\[
				\begin{aligned}
					L &= 2\sqrt{r^{2}-d^{2}} \\
					&= 2\sqrt{25-9} \\
					&= 2\sqrt{16} \\
					&= 8
			\end{aligned}
			\]
			\par\textbf{理由：} 直接应用弦长公式。
	\end{description}

	\textbf{\textcolor{CExample}{关键结论：}}
	\textbf{弦长为 $8$。}

	\medskip
	\textbf{\textcolor{CExample}{例 3（综合应用）}}

	\textbf{\textcolor{CExample}{题目：}}
	过点 $P(0,5)$ 作直线交圆 $x^{2}+y^{2}=25$ 于两点，若截得的弦长为 $6$，求直线的方程。

	\textbf{\textcolor{CExample}{解题步骤：}}
	\begin{description}[style=nextline, leftmargin=2.8em, labelsep=0.5em, itemsep=0.45em]
		\item[\textbf{第一步（反求弦心距）：}] 由 $L=6$，$r=5$，得
			\[
				\begin{aligned}
					d &= \sqrt{r^{2}-\left(\frac{L}{2}\right)^{2}} \\
					&= \sqrt{25-9} \\
					&= 4
			\end{aligned}
			\]
			\par\textbf{理由：} 弦长公式逆用，求弦心距。
		\item[\textbf{第二步（设直线方程）：}] 设直线斜率为 $k$，则方程为 $y=kx+5$，即 $kx-y+5=0$。
			\par\textbf{理由：} 点斜式，斜率待定。
		\item[\textbf{第三步（由距离求斜率）：}] 圆心 $(0,0)$ 到直线距离等于 $4$，即 $\frac{|5|}{\sqrt{k^{2}+1}}=4$，解得 $k^{2}+1=\frac{25}{16}$，$k^{2}=\frac{9}{16}$，$k=\pm\frac{3}{4}$。
			\par\textbf{理由：} 点到直线距离公式。
		\item[\textbf{第四步（写出方程并检验）：}] 当 $k=\frac{3}{4}$ 时，$3x-4y+20=0$；当 $k=-\frac{3}{4}$ 时，$3x+4y-20=0$。经检验斜率不存在时弦长为 $10\neq6$，舍去。
			\par\textbf{理由：} 代入斜率得直线，并检查特殊情况。
	\end{description}

	\textbf{\textcolor{CExample}{关键结论：}}
	\textbf{直线方程为 $3x-4y+20=0$ 或 $3x+4y-20=0$。}
\end{examplebox}
